\documentclass[12pt,letterpaper]{article}
\usepackage[utf8]{inputenc}
\usepackage[french]{babel}
\usepackage{amsmath, amssymb, amsfonts}
\usepackage{graphicx}
\usepackage{hyperref}
\usepackage{enumitem}
\usepackage[margin=2.5cm]{geometry}
\usepackage[T1]{fontenc}

\title{Leçon 2 : Virgule flottante et erreurs numériques}
\author{Analyse Numérique}
\date{Séance en classe}

\begin{document}

\maketitle

\begin{center}
\textbf{Objectif :} Comprendre comment un ordinateur stocke les nombres réels et pourquoi cela peut provoquer des catastrophes dans les calculs scientifiques.
\end{center}

\section*{1. Le standard IEEE 754 (virgule flottante)}

Dans un ordinateur, un nombre réel est représenté en binaire avec trois composantes :

\[
\text{Nombre} = (-1)^s \times (1 + \text{mantisse}) \times 2^{\text{exposant} - \text{biais}}
\]

\subsection*{Les trois composantes}
\begin{itemize}
    \item \textbf{Le bit de signe ($s$) :} 0 pour positif, 1 pour négatif.
    \item \textbf{L'exposant ($e$) :} un entier décalé par un biais.
    \item \textbf{La mantisse ($m$) :} la partie fractionnaire (toujours de la forme $1,m$).
\end{itemize}

\subsection*{Les deux formats les plus courants}
\begin{center}
\begin{tabular}{|c|c|c|c|}
\hline
\textbf{Format} & \textbf{Signe} & \textbf{Exposant} & \textbf{Mantisse} \\
\hline
Simple précision (32 bits)  & 1 bit & 8 bits & 23 bits \\
Double précision (64 bits)  & 1 bit & 11 bits & 52 bits \\
\hline
\end{tabular}
\end{center}

\subsection*{Exemple : représentation de $0.3$ en binaire}
Nous avons vu que $0.3_{10} = 0.01001100110011\ldots_2$ (périodique). En notation scientifique normalisée :

\[
0.3_{10} = 1.001100110011\ldots_2 \times 2^{-2}
\]

En double précision, l'ordinateur ne peut stocker que 52 bits de la mantisse. Il tronque ou arrondit le reste. Ce petit écart est la source de l'erreur.

\subsection*{Pourquoi $0.2 + 0.3 \neq 0.5$ parfois ?}
En réalité, $0.2$ et $0.3$ sont arrondis dès leur entrée dans l'ordinateur. Leur somme est donc la somme de deux approximations. Parfois, les erreurs s'annulent ; parfois, elles s'additionnent.

\section*{2. Effet des erreurs sur les méthodes numériques}

\subsection*{Exemple 1 : La soustraction de deux nombres très proches (phénomène de cancellation)}

Soit à calculer $f(x) = \sqrt{x+1} - \sqrt{x}$ pour $x$ grand.

Pour $x = 10^6$ :
\[
\sqrt{10^6 + 1} - \sqrt{10^6} = 1000.0005 - 1000 = 0.0005
\]

En virgule flottante, $\sqrt{10^6+1}$ et $\sqrt{10^6}$ sont presque égaux. La soustraction fait perdre des chiffres significatifs. Il est préférable de réécrire :
\[
\sqrt{x+1} - \sqrt{x} = \frac{1}{\sqrt{x+1} + \sqrt{x}}
\]

\textbf{Morale :} Les soustractions de nombres presque égaux sont dangereuses.

\subsection*{Exemple 2 : La méthode de Horner vs l'évaluation directe}

Soit le polynôme $P(x) = (x-1)^{10}$ développé :
\[
P(x) = x^{10} - 10x^9 + 45x^8 - 120x^7 + 210x^6 - 252x^5 + 210x^4 - 120x^3 + 45x^2 - 10x + 1
\]

Si on évalue ce polynôme en $x = 0.999$ avec la forme développée, les grands coefficients ($-252$, $210$, etc.) créent des erreurs d'arrondi. En utilisant la forme factorisée $(x-1)^{10}$, on obtient un résultat beaucoup plus précis.

\textbf{Morale :} La forme d'évaluation d'une expression influence la stabilité numérique.

\subsection*{Exemple 3 : Somme de séries alternées}

Soit la série $S = \sum_{n=1}^{\infty} \frac{(-1)^{n+1}}{n}$. En théorie, $S = \ln 2 \approx 0.693147$.

Si on somme les termes en commençant par les grands et en ajoutant les petits, on peut accumuler des erreurs. Si on commence par les petits, l'erreur est moindre. L'ordre des opérations compte !

\subsection*{Exemple 4 : Équation quadratique classique}

Soit l'équation $x^2 - 2x + 0.9999 = 0$. Les racines sont $1 \pm 0.01$, soit $1.01$ et $0.99$.

La formule classique :
\[
x = \frac{-b \pm \sqrt{b^2 - 4ac}}{2a}
\]
si $b$ est très grand par rapport à $4ac$, peut donner une soustraction de deux nombres très proches pour l'une des racines.

On peut utiliser la forme alternative :
\[
x_1 = \frac{2c}{-b - \sqrt{b^2 - 4ac}}, \quad x_2 = \frac{-b - \sqrt{b^2 - 4ac}}{2a}
\]
pour éviter la cancellation.

\section*{3. Exercices pour la séance}

\begin{enumerate}
    \item Calculez en Python / Julia les trois expressions suivantes et expliquez les résultats :
    \begin{verbatim}
    0.1 + 0.2 == 0.3
    0.1 + 0.2 - 0.3
    (0.1 + 0.2) - 0.3
    \end{verbatim}
    
    \item Évaluez le polynôme $P(x) = (x-1)^6$ en $x = 1.001$ en utilisant la forme développée et la forme factorisée. Comparez les résultats.

    \item Calculez $\sum_{n=1}^{1000} \frac{1}{n^2}$ en ajoutant les termes de $1$ à $1000$ et de $1000$ à $1$. Comparez avec $\pi^2/6 \approx 1.644934$.
\end{enumerate}

\section*{4. Résumé à retenir}

\begin{center}
\fbox{\begin{minipage}{0.9\textwidth}
\begin{itemize}
    \item Les nombres ne sont pas stockés exactement.
    \item Les erreurs s'accumulent.
    \item La soustraction de termes proches est dangereuse.
    \item L'ordre des opérations et la forme de l'expression influencent la précision.
    \item Un bon algorithme doit être \textbf{stable numériquement}.
\end{itemize}
\end{minipage}}
\end{center}

\vspace{1cm}
\textit{“En analyse numérique, il ne s'agit pas de trouver la bonne réponse, mais de trouver une réponse qui ne soit pas trop fausse.”}

\end{document}